% !TEX root = ../notes.tex

\documentclass[../notes.tex]{subfiles}

\begin{document}

\section{April 22}
Today, we discuss the Cartan decomposition.

\subsection{The Cartan Decomposition}
Here is our statement.
\begin{theorem}
	Fix a real form $\mf g_\theta$ of a semisimple complex Lie algebra $\mf g$, and let $\mf g=\mf p\oplus\mf k$ be the decomposition into $\theta$-eigenspaces. We then set $\mf g_\theta=\mf k^c\oplus\mf p_\theta$. Let $G_\theta\subseteq G_{\mathrm{ad}}$ be the invariants, and we define $K^c\subseteq G_\theta$ and $A\subseteq G_\theta$ to be the subgroup corresponding $\mf k^c$ and to a maximal abelian subspace $\mf a\subseteq\mf p_\theta$, respectively. Then
	\[G_\theta=K^cAK^c.\]
\end{theorem}
\begin{proof}
	The polar decomposition (\Cref{thm:polar-decomp}) grants $G_\theta=K^cP_\theta$, so it is enough to show that every element in $P_\theta$ can be found in $K^cAK^c$. We will show that every element in $P_\theta$ can be conjugated by $K^c$ into $A$. Well, choose $y\in P_\theta$, which we may write as $\log Y$ for some unique $Y\in\mf p_\theta$. By \Cref{prop:maximal-abelian-subspace}, we find that there is $k\in K^c$ with $\op{Ad}_ky\in\mf a$. Taking the exponential shows that $\op{Ad}_kY\in A$, so we are done.
\end{proof}
\begin{remark}
	There is a generalization to reductive groups, which is no harder to prove.
\end{remark}
\begin{remark}
	The decomposition is not unique.
\end{remark}
\begin{example}
	Take $G_\theta=\op{GL}_n(\CC)$. Then $K=\op U(n)$ and $A$ can be taken to be the subgroup of diagonal operators with positive eigenvalues. Thus, we are saying that any $g\in\op{GL}_n(\CC)$ can be decomposed as $u_1du_2$ where $u_1,u_2\in\op U(n)$ and $d\in A$.
\end{example}
\begin{example} \label{ex:cartan-gln-r}
	Take $G_\theta=\op{GL}_n(\RR)$. Then $K=\op O(n)$ and $A$ can be taken to be the subgroup of diagonal operators with positive eigenvalues. Thus, we are saying that any $g\in\op{GL}_n(\RR)$ can be decomposed as $o_1do_2$, where $o_2,o_2\in\op O(n)$ and $d\in A$.
\end{example}

\subsection{Maximal Compact Subgroups}
We are now able to prove the most important fact about the $K^c$s: they are maximal compact subgroups, and with this property, they are unique up to conjugation.
\begin{theorem}[Cartan] \label{thm:cartan-compact}
	Let $G_\theta$ be a real form of a connected semisimple complex Lie group $G$. Then any compact subgroup $L\subseteq G^\theta$ is conjugate (by $G_\theta$) to a subgroup of $K^c$.
\end{theorem}
Note that this gives the following promised corollary.
\begin{corollary}
	Let $G_\theta$ be a real form of a connected semisimple complex Lie group $G$. Then all maximal compact subgroups in $G_\theta$ are conjugate to $K^c$.
\end{corollary}
\begin{proof}
	This follows from \Cref{thm:cartan-compact}: any maximal compact subgroup $L$ is conjugate to a subgroup of $K^c$, but the conjugate must remain maximal, so $L$ must be conjugate to $K^c$ on the nose. (Note that $K^c$ is compact because it is closed in $G^c$.)
\end{proof}
We now prove \Cref{thm:cartan-compact}, following Mostow. We start by noticing that it is not actually too hard to show the maximality of $K^c$.
\begin{theorem}
	Let $G_\theta$ be a real form of a connected semisimple complex Lie group $G$. Then $K^c$ is a maximal compact subgroup.
\end{theorem}
\begin{proof}
	Suppose that $L$ is a compact subgroup containing $K^c$. By the polar decomposition, any $k\in K$ can be decomposed into $K^c$ and $P_\theta$, but then the element in $P_\theta$ must still be in $L$. In other words, $L=K^c(P_\theta\cap L)$.

	We now aim to show that $P_\theta\cap L$ is trivial, which will complete the proof. Well, choose some $y\in P_\theta\cap L$. Then the eigenvalues of $y$ act on $\mf g$ by positive Hermitian operators. Thus, if any of the eigenvalues are not one, then at least one of the eigenvalues is bigger than one, so the eigenvalues of $\{y^n\}_{n\ge1}$ escape to $\infty$, so there is no convergent subsequence of the image of these elements in $\op{GL}(\mf g)$, which contradicts the compactness of $L$! Thus, $y=1$, and we are done.
\end{proof}
Let's now turn to the rest of the proof of \Cref{thm:cartan-compact}
\begin{proof}[Proof of \Cref{thm:cartan-compact}]
	We proceed in steps. Choose a compact subgroup $L\subseteq G_\theta$, and we want to show that it is conjugate to a subgroup of $K^c$ by $G_\theta$. We may as well show that $L$ is conjugate by $P_\theta$ by \Cref{thm:polar-decomp}. The idea is to construct an $L$-invariant function $f\colon P_\theta\to\RR$ with a unique minimum $\exp(Y)\in P_\theta$, and this minimum will encode the correct conjugate. Indeed, suppose we have such a function $f$. Then because $f$ is $L$-invariant, its unique minimum $y\in P_\theta$ must be an $L$-invariant point. Thus, for all $h\in L$, we have $h\exp(Y)=\exp(Y)\overline h$. This rearranges into
	\[\exp(-Y/2)h\exp(Y/2)=\exp(Y/2)\overline h\exp(-Y/2),\]
	which implies that $\exp(-Y/2)h\exp(Y/2)$ is conjugate to itself, meaning that it is in $K$.

	It now remains to hunt for this function $f$.
	\begin{enumerate}
		\setcounter{enumi}{0}
		\item We set some notation. Let $B$ be the Killing form on $\mf g$, restricted to $\mf g_\theta$; thus, it is negative-definite on $\mf k^c$ and positive-definite on $\mf p_\theta$. We then define $B_\theta$ by
		\[B_\theta(x,y)\coloneqq-B(x,\theta y),\]
		which is positive-definite on $\mf g_\theta$ and $K^c$-invariant (because $K^c$ is $\theta$-invariant). Now, for an operator $A$ on $\mf g_\theta$, we let $A^\dagger$ be the adjoint under $B_\theta$. Note the following special cases.
		\begin{itemize}
			\item For $g\in K^c$, the operator $\op{Ad}_g$ is orthogonal because it preserves $B_\theta$.
			\item For $g\in P_\theta$, the operator $\op{Ad}_g$ is self-adjoint and has positive eigenvalues.
		\end{itemize}
		Accordingly, for $g=kp$ with $k\in K^c$ and $p\in P_\theta$ (given by \Cref{thm:polar-decomp}), we find that
		\begin{align*}
			\mathrm{Ad}_g^\dagger &= \mathrm{Ad}_p^\dagger\mathrm{Ad}_k^\dagger \\
			&= \mathrm{Ad}_p\mathrm{Ad}_{k^{-1}} \\
			&= \mathrm{Ad}_{\overline g}^{-1}
		\end{align*}
		by construction of the complex conjugate.

		\item We define the function $f$. To start, we define
		\[S\coloneqq\int_L\mathrm{Ad}_h^\dagger\mathrm{Ad}_h\,dh,\]
		which is some endomorphism of $\mf g_\theta$. Note that $(\cdot)^\dagger$ is linear on operators, so $S^\dagger=S$. Additionally, $S$ is positive because it is an average of positive operators. Thus, $S$ can be given an orthonormal eigenbasis $\{v_i\}$ with positive eigenvalues $\{\lambda_i\}$. We let $\lambda$ be the minimal eigenvalue and then define $f\colon P_\theta\to\RR$ by
		\[f(x)\coloneqq\tr(\mathrm{Ad}_x\circ S).\]
		Note that taking the trace can be computed as
		\[f(x)=\sum_i\lambda_iB_\theta(\op{Ad}_xv_i,v_i).\]
		Because $\op{Ad}_x$ is positive and self-adjoint, this is then bounded below by $\lambda\tr\mathrm{Ad}_x$. While we're here, we also note that $f$ is $L$-invariant because $G_\theta$ acts on $P_\theta$ by the twisted adjoint action $g\cdot x=gx\overline g^{-1}$.
		
		\item We show that $f$ has a minimum $y\in P_\theta$. Well, for any $R>0$, note that the set
		\[\{A\in M_n(\RR):A=A^\intercal,\det A=1,\tr A\le R\}\]
		is compact. Indeed, up to the orthogonal sets $\op O_n(\RR)$, any such matrix $A$ can be diagonalized with positive entries by \Cref{ex:cartan-gln-r}, but then the diagonal entries have positive coefficients and are bounded above by $R$ (with product $1$). Thus, the collection of relevant diagonal matrices is compact, so the entire set is compact.

		Thus, for large $R$, we conclude that the set of $X$ with $f(X)\le R$ is a compact set because they are some self-adjoint operators with bounded trace (and determinant $1$ on $\mf g$). We conclude that $f$ attains a minimum on $P_\theta$.

		\item As an intermission, we discuss a small piece of linear algebra: given symmetric real matrices $a$ and $M$ with $a$ nonzero and $M$ positive-definite. Then we claim that the function
		\[\varphi(t)\coloneqq\tr(\exp(ta)M)\]
		is strictly convex. Indeed, by changing basis, we may assume that $a$ is a diagonal matrix with diagonal entries $\{a_i\}_i$. Then
		\[\tr(\exp(ta))=\sum_i\exp(ta_i)M_{ii}.\]
		Because $M$ is positive-definite, the entries $M_{ii}$ are positive, so this is a positive linear combination of strictly convex functions and hence strictly convex.

		\item We now show that $f$ admits a unique minimum $y\in P_\theta$. Suppose that we have another minimum $z\in P_\theta$ of $f$. Set $Y\coloneqq\log y$ and $Z\coloneqq\log z$, and define $g\coloneqq\exp(Z/2)\exp(-Y/2)$. We would like to show that $g=1$, which will imply that $Z/2=Y/2$ by taking logarithms, which then implies $y=z$ by taking exponentials.

		To show that $g=1$, we use the polar decomposition: by \Cref{thm:polar-decomp}, we have $k\in K^c$ and $X\in\mf p_\theta$ with
		\[\exp\left(\frac Z2\right)\exp\left(-\frac Y2\right)=k\exp\left(\frac X2\right).\]
		Thus,
		\[\exp\left(\frac X2\right)=k^{-1}\exp\left(\frac Z2\right)\exp\left(-\frac Y2\right).\]
		Because $\exp(X/2)$ is self-adjoint, we also see that
		\[\exp\left(\frac X2\right)=\exp\left(-\frac Y2\right)\exp\left(\frac Z2\right)k,\]
		so multiplying gives
		\[\exp(X)=\exp\left(-\frac Y2\right)\exp(Z)\exp\left(-\frac Y2\right),\]
		which rearranges to
		\[\exp(Z)=\exp\left(-\frac Y2\right)\exp(X)\exp\left(-\frac Y2\right).\]
		We now consider the function
		\[g(t)\coloneqq f\left(\exp\left(-\frac Y2\right)\exp(tX)\exp\left(-\frac Y2\right)\right).\]
		By construction, this function attains minima at $t=0$ and $t=1$. On the other hand,
		\[g(t)=\tr\left(\mathrm{Ad}_{\exp(-Y/2)}\circ\mathrm{Ad}_{\exp(tX)}\circ\mathrm{Ad}_{\exp(-Y/2)}\circ S\right),\]
		which is
		\[g(t)=\tr\left(\mathrm{Ad}_{\exp(tX)}\circ\left(\mathrm{Ad}_{\exp(-Y/2)}\circ S\circ\op{Ad}_{\exp(Y/2)}\right)\right)\]
		takes the form of the functions in the previous step (recall $S$ is self-adjoint and positive), so it is strictly convex (with a unique minimum) as soon as $X$ is nonzero. Thus, we instead must have $X=0$, so $g\in K^c$. But then $\exp(Z/2)=k\exp(Y/2)$, so the polar decomposition forces $k=1$, so we are done.
		\qedhere
	\end{enumerate}
\end{proof}

\subsection{Cartan Subalgebras}
We now take a moment to discuss Cartan subalgebras.
\begin{definition}[Cartan]
	A \textit{Cartan subalgebra} of a real form $\mf g_\theta$ is a maximal abelian subalgebra consisting of semisimple elements.
\end{definition}
\begin{remark}
	By construction, one finds that $\mf h$ is a Cartan subalgebra if and only if $\mf h_\CC$ is a Cartan subalgebra of $\mf g_\theta$. The main point is that any maximal commutative subalgebra $\mf h$ with semisimple elements, and we want to show that $\mf h_\CC$ continues to be maximal. Indeed, if $X\in\mf g$ commutes with $\mf h_\CC$, then so does $\overline X$, so the real part $(X+\overline X)/2$ and the imaginary part $(X-i\overline X)/2i$ are both in $\mf h$, so $X\in\mf h_\CC$.
\end{remark}
\begin{example}
	For $\mf g=\mf{sl}_n(\RR)$, choose some nonnegative integer $m\le n/2$. Then there is a Cartan subalgebra $\mf h_m$ consisting of the block-diagonal matrices
	\[\left\{\op{diag}\left(\begin{bmatrix}
		a_1 & b_1 \\ -b_1 & a_1
	\end{bmatrix},\ldots,\begin{bmatrix}
		a_m & b_m \\ -b_m & a_m
	\end{bmatrix},c_1,\ldots,c_{n-2m}\right):2(a_1+\cdots+a_m)+(c_1+\cdots+c_{n-2m})=0\right\}.\]
	On the homework, we will show that these Cartan subalgebras exhaust all of them up to isomorphism.
\end{example}
Let's give some invariants to classify Cartan subalgebras.
\begin{definition}[split]
	Fix a real form $\mf g_\theta$. A semisimple element $X\in\mf g_\theta$ is \textit{split} if and only if the adjoint action of $X$ on $\mf g_\theta$ has real eigenvalues. A commutative subalgebra $\mf h$ of $\mf g_\theta$ is \textit{split} if and only if all elements of $\mf h$ are split; for any commutative subalgebra $\mf h$, we let $\mf h^{\mathrm{spl}}$ be the subspace of split elements, and we define $s(\mf h)\coloneqq\dim\mf h^{\mathrm{spl}}$.
\end{definition}
\begin{remark}
	The commutativity implies that $\mf h^{\mathrm{spl}}$ is in fact a subspace.
\end{remark}
\begin{remark}
	Note that $\mf h$ is split if and only if $\mf h^{\mathrm{spl}}=\mf h$, which is equivalent to $s(\mf h)=\dim\mf h$.
\end{remark}
\begin{remark}
	One can check that the signature of the Killing form restricted to $\mf h$ is $(\dim\mf h-s(\mf h),s(\mf h))$.
\end{remark}
\begin{remark}
	One can check that $\exp(\mf h)\cong\left(S^1\right)^{\dim\mf h-s(\mf h)}\times\left(\RR^+\right)^{s(\mf h)}$.
\end{remark}
\begin{remark}
	It turns out that $\mf g_\theta$ is split if and only if it admits a split Cartan subalgebra. On one hand, being the split form allows the usual $\mf h$, which is split. Conversely, having a split Cartan means that its root decomposition can be taken to have real eigenvalues everywhere (by definition) and thus fully defined over $\RR$, which allows us to present $\mf g_\theta$ as the split real form.
\end{remark}
It is interesting to determine what values of $s(\mf h)$ are possible.
\begin{definition}
	Fix a real form $\mf g_\theta$. Then $\mf h$ is \textit{maximally split} if and only if $s(\mf h)$ is as large as possible. Similarly, $\mf h$ is \textit{maximally compact} if and only if $s(\mf h)$ is as small as possible.
\end{definition}
\begin{example}
	For $\mf{sl}_n(\RR)$, the diagonal torus $\mf h_0$ is maximally split, and the diagonal torus $\mf h_m$ is maximally compact.
\end{example}
\begin{remark}
	Any compact Cartan subalgebra has $s(\mf h)$ equal to the rank, which is automatically maximally compact.
\end{remark}
We will show the following result next time.
\begin{proposition}
	Fix a real form $\mf g_\theta$, and let $\mf h$ be a $\theta$-stable Cartan subalgebra, and let $\mf g_\theta=\mf k^c\oplus\mf p_\theta$ be the polar decomposition.
	\begin{listalph}
		\item Then $\mf h$ is maximally split if and only if $\mf h\cap\mf p_\theta$ is a maximal abelian subspace of $\mf p_\theta$
		\item Then $\mf h$ is maximally compact if and only if $\mf h\cap\mf k^c$ is a Cartan subalgebra.
		\item Any two maximally split Cartan subalgebras are conjugate.
		\item Any two maximally compact Cartan subalgebras are conjugate.
		\item Any Cartan subalgebra is conjugate to a $\theta$-stable one (by $G_\theta$).
	\end{listalph}
\end{proposition}
\begin{proof}
	Omitted. It's present in the notes.
\end{proof}

\end{document}